% Part: lambda-calculus
% Chapter: church-rosser
% Section: beta-eta-reduction

\documentclass[../../../include/open-logic-section]{subfiles}

\begin{document}

\olfileid[pt]{lam}{cr}{be}

\olsection{$\beta\eta$-redução}

A propriedade de Church--Rosser vale para a $\beta\eta$-redução
($\bered$).

\begin{lem}\ollabel{lem:one-par}
  Se $M \beredone M'$, então $M \beredpar M'$.
\end{lem}

\begin{proof} 
  Por indução na !!{derivation} de $M \beredone M'$. Se $M \bredone M'$
  por conversão-$\eta$ (isto é, \olref[syn][eta]{defn:beredone}), usamos
  \olref[pbe]{thm:refl}. Os outros casos são como em
  \olref[b]{lem:one-par}.
\end{proof}


\begin{lem}\ollabel{lem:par-red}
  Se $M \beredpar M'$, então $M \bered M'$.
\end{lem}

\begin{proof} Indução na !!{derivation} de $M \beredpar M'$.

  Se a última regra é \olref[pbe]{defn:beredpar5}, então $M$ é
  $\lambd[x][Nx]$ e $M'$ é $N'$ para alguns $x$, $N$, $N'$, onde
  $x \notin FV(N)$ e $N \beredpar N'$. Assim, primeiro podemos reduzir
  $\lambd[x][Nx]$ a $N$ por conversão-$\eta$, seguida da série de passos
  $\beredone$ que mostra que $N \bered N'$, o que vale pela hipótese de
  indução.
\end{proof}


\begin{lem}\ollabel{lem:str}
  $\bered$ é a menor relação transitiva que contém $\beredpar$.
\end{lem}

\begin{proof}
  Como em \olref[b]{lem:str}.
\end{proof}

\begin{thm}\ollabel{thm:cr}
  $\bered$ satisfaz a propriedade de Church--Rosser.
\end{thm}

\begin{proof}
  Por \olref[dap]{thm:str}, \olref[pbe]{thm:cr} e \olref{lem:str}.
\end{proof}
\end{document}

